% ID: SEL-D0062-OA-E04
% ESTADO: cotejada. Segunda pasada visual y borrador compilado revisados.
% Fuente: Selectividad/sel_2010_matematicas_aplicadas/mat_apli/mat_apli_junio/mat_apli_junio_ra/EXAMEN-ANDALUCÍA-4.pdf
% SHA-256: 4e10e40d60aebdab7570fa77b9ce953d7bf6c9a4c8869a5d3f9b7ca3c907c154
% Página: 1; opción A; ejercicio original 4.

Se desea estimar la proporción de votantes a un determinado partido político mediante una muestra aleatoria.
\begin{enumerate}[label=\alph*)]
\item \textbf{(1.25 puntos)} Si de una muestra de 500 personas 200 dicen que lo votan, calcule con un nivel de confianza del $97\%$ un intervalo para la proporción de votantes a ese partido en la población.
\item \textbf{(1.25 puntos)} Si la proporción de votantes en otra muestra ha sido $0.2$ y el error cometido en la estimación ha sido inferior a $0.05$, con un nivel de confianza del $99\%$, calcule el tamaño mínimo de dicha muestra.
\end{enumerate}
